\chapter{The OVO System}
\label{ch:preliminaries}

% PRESUPUESTO: ~5 pp
% Capítulo de EXPOSICIÓN: explicar OVO a quien no lo conoce, inclinado hacia lo que
% la tesis ataca, pero SIN juzgar todavía (eso es Ch5). Termina en limitaciones =
% trampolín al resto de la tesis. El setup + baseline se han movido a Ch4 (Benchmark).

OVO~\parencite{berrielmartins2024ovo} es un sistema de mapeo semántico 3D
online de vocabulario abierto sobre el que se construye el conjunto
de esta tesis. Cualquier mejora sobre él exige primero comprender qué produce el sistema y
cómo se articulan sus componentes; sobre esa base es posible identificar qué elementos
inciden en el objeto de estudio de este trabajo y razonar por dónde conviene intervenir.

El capítulo recorre esos fundamentos en orden: la representación 3D que genera, la
arquitectura con que acopla un backbone de SLAM a la segmentación semántica de instancias,
el mecanismo de fusión de instancias en loop closure sobre el que actúan las contribuciones
posteriores y, para cerrar, las limitaciones que abren el resto de la tesis.

\section{System overview}
\label{sec:ovo_overview}

OVO es un sistema que, a partir de un flujo de keyframes RGB-D con sus poses, produce una
\textbf{reconstrucción densa} de la escena: una nube de puntos 3D en la que cada objeto se
recupera como una \emph{instancia 3D} distinta, descrita por un vector semántico de vocabulario
abierto. A diferencia de los métodos de vocabulario cerrado, la instancia no se clasifica en un
conjunto fijo de categorías: se le asocia un descriptor calculado con un modelo visión-lenguaje
(CLIP,~\cite{radford2021clip}), que proyecta imágenes y texto a un mismo espacio. Así la escena
puede consultarse con lenguaje natural, por similitud de coseno entre el texto de la consulta y
los descriptores de las instancias (Figura~\ref{fig:ovo_queries}).

% FIGURA: Fig. 3 del paper (queries open-vocab). images/ovo_queries.png
\begin{figure}[h!]
  \centering
  \includegraphics[width=\textwidth]{ovo_queries}
  \caption{Open-vocabulary queries on an OVO map: each natural-language query highlights its
    matching instances. Adapted from~\cite{berrielmartins2024ovo}.}
  \label{fig:ovo_queries}
\end{figure}

OVO construye esa representación de forma \emph{online} e incremental: en lugar de reprocesar la
escena entera, integra cada nuevo keyframe sobre el mapa existente, que crece a medida que la
cámara recorre el entorno (Figura~\ref{fig:ovo_semantic_map}).

% FIGURA: Fig. 1 del paper (mapa semántico creciendo). images/ovo_semantic_map.png
\begin{figure}[h!]
  \centering
  \includegraphics[width=0.6\textwidth]{ovo_semantic_map}
  \caption{Output of OVO on a Replica sequence: an open-vocabulary semantic map built
    incrementally as the camera explores the scene. Adapted from~\cite{berrielmartins2024ovo}.}
  \label{fig:ovo_semantic_map}
\end{figure}

Para sostener ese ritmo, OVO adopta el principio de parallel tracking and
mapping~\parencite{klein2007ptam}, habitual en SLAM visual: mantener un hilo rápido siguiendo la
cámara y desacoplar de él el trabajo más pesado. OVO aplica ese mismo principio para separar dos trabajos que
ejecuta de forma asíncrona:

\begin{itemize}
  \item \textbf{Reconstrucción geométrica.} Retroproyecta a 3D la profundidad de cada keyframe
    usando las poses $T_n$ que estima el backbone de SLAM, levantando la nube de puntos densa sobre
    la que actúa la semántica. De dónde venga la pose, un SLAM real o una referencia ground-truth,
    le es indiferente: la geometría es solo el terreno sobre el que trabaja.
  \item \textbf{Identificación semántica.} Combina tres ingredientes: las máscaras 2D que un
    segmentador (SAM~2.1,~\cite{ravi2024sam2}) extrae de cada keyframe, la reproyección de la nube
    al plano de imagen y un descriptor de vocabulario abierto (CLIP). Reproyectando la nube bajo
    cada máscara, un mecanismo de asociación decide si crea una instancia nueva o propaga una ya
    existente a esos puntos (Sección~\ref{sec:ovo_segment_mapper}). A cada instancia se le asigna
    además ese descriptor, que codifica su semántica y permite consultarla con lenguaje natural
    (Sección~\ref{sec:ovo_descriptors}).
\end{itemize}

Ambos escriben sobre un mismo mapa $\mathcal{M} = \{\mathcal{T}, \mathcal{P}, \mathcal{S}\}$: las
poses de los keyframes $\mathcal{T}$, la nube de puntos densa $\mathcal{P}$ y el conjunto de
instancias 3D $\mathcal{S}$. Cada punto $P = ([x\,y\,z]^\top, l_p)$ reúne las dos aportaciones: las
coordenadas 3D que fija la geometría y una \textbf{etiqueta discreta} $l_p \in \{-1, 0, 1, \dots,
q\}$ que apunta a su instancia y que rellena la semántica ($l_p = -1$ mientras está sin asignar).

\section{Geometric reconstruction}
\label{sec:ovo_geometric}

De los tres componentes del mapa (Sección~\ref{sec:ovo_overview}), las poses $\mathcal{T}$ las
estima el SLAM y las instancias $\mathcal{S}$ son cosa de la parte semántica
(Sección~\ref{sec:ovo_semantic}). La reconstrucción geométrica se ocupa del tercero, la nube de
puntos densa $\mathcal{P}$: la construye a partir de las poses, sin estimarlas él, para levantar y
mantener la geometría sobre la que actúa el resto del sistema. Lo hace en tres operaciones:

\begin{itemize}
  \item \textbf{Construcción de la nube.} Por cada keyframe $k_n$, OVO añade a $\mathcal{P}$ los
    puntos 3D de la escena. Si el backbone de SLAM ya estima una nube densa, OVO concatena esos
    puntos directamente; si no, los obtiene retroproyectando a 3D la profundidad del keyframe con
    la pose $T_n$. En la práctica, un SLAM denso (Gaussian-SLAM) aporta la nube, mientras que con
    un SLAM sparse (ORB-SLAM2) o con poses ground-truth OVO recurre a la retroproyección. Cada
    punto nace con $l_p = -1$, es decir, sin semántica asignada todavía.
  \item \textbf{Coherencia geométrica.} Para que $\mathcal{P}$ no crezca sin control al revisitar
    una zona, OVO no añade puntos donde ya los hay. Proyecta los existentes al keyframe y descarta
    cada píxel donde un punto ya proyectado cae en su vecindad \emph{y} coincide en profundidad con
    la que mide el sensor, dentro de un umbral; esa coincidencia distingue una superficie revista
    de otra oculta detrás. Los píxeles restantes, superficie aún no vista, sí generan puntos.
  \item \textbf{Corrección con el loop closure.} El SLAM acumula deriva, así que un mismo objeto
    visto en dos pasadas puede quedar \textbf{duplicado} en dos instancias separadas. Al cerrar un
    loop, el SLAM corrige la trayectoria y OVO re-transforma la nube arrastrando la etiqueta $l_p$
    de cada punto: la geometría se corrige sin perder la semántica, pero los dos duplicados quedan
    ahora \textbf{superpuestos}. Deshacerlos es lo que ataca el mecanismo de fusión de instancias,
    disparado en ese instante (Sección~\ref{sec:ovo_fusion_lc}).
\end{itemize}

\section{Semantic identification}
\label{sec:ovo_semantic}

Sobre la nube que produce la reconstrucción geométrica, la identificación semántica asigna cada punto
a una instancia y describe cada instancia para poder consultarla. Sus tres piezas se detallan a
continuación: cómo se asocian y siguen las instancias frame a frame, cómo se calcula su descriptor
y cómo se fusionan cuando el SLAM cierra un loop.

\subsection{3D instance segmentation and tracking}
\label{sec:ovo_segment_mapper}

La reconstrucción geométrica entrega una nube de puntos sin etiquetar ($l_p = -1$). Poblarla de
instancias, creándolas y siguiéndolas a medida que la cámara avanza, es la tarea de esta primera
pieza: aquí nace $\mathcal{S}$, el tercer componente del mapa (Sección~\ref{sec:ovo_overview}).

\paragraph{La instancia.} Una instancia 3D queda descrita por $S = (\mathbf{d}, \kappa)$: un
descriptor semántico $\mathbf{d}$ y una lista $\kappa$ de los keyframes desde los que mejor se
observa. El puente entre la nube y las instancias es la etiqueta \emph{por punto} $l_p$
(Sección~\ref{sec:ovo_overview}): enlaza cada punto con su instancia, es la unidad más fina de la
representación y la variable sobre la que actúan las contribuciones posteriores.

\paragraph{El bucle de asociación.} El 3D segment mapper (Algoritmo~\ref{alg:segment_mapper}) es el
mecanismo responsable de la creación y el tracking de instancias. Por cada keyframe $k_n$ decide qué
máscaras 2D corresponden a instancias ya vistas y cuáles son objetos nuevos, mediante una
\textbf{asociación por proyección y votación}. SAM trocea $k_n$ en máscaras
$\tilde{\mathcal{S}}_n$; los puntos del mapa visibles bajo la pose $T_n$ se proyectan al plano de
imagen (descartando ocluidos por profundidad), arrastrando cada uno su etiqueta $l_p$. Para cada
máscara $s$ se toma la \emph{moda} de esas etiquetas, $z_l := \arg\max_{l_p}
(\tilde{\mathcal{P}}_n \cap s)$, con sus votos $v$. Si $v \le \epsilon$ la máscara se descarta por
falta de evidencia; si $z_l = -1$ es un \textbf{objeto nuevo} (se crea $S_{q+1}$ con descriptor
vacío); y si $z_l$ es una instancia existente, la máscara \textbf{es esa instancia vista de nuevo}
y se le añade $k_n$ como vista. Finalmente la etiqueta ganadora se propaga a los puntos
\emph{sin asignar} bajo la máscara y a la propia máscara, y $(k_n, \tilde{\mathcal{S}}_n)$ se
encola en $\mathcal{Q}$ para el descriptor.

% ALGORITMO: Algorithm 1 del paper (3D Segment Mapper).
\begin{algorithm}[h!]
  \refstepcounter{algorithm}\label{alg:segment_mapper}%
  \hrule height 0.8pt\vspace{3pt}
  \noindent\textbf{Algorithm~\thealgorithm}\quad 3D Segment Mapper\par
  \vspace{2pt}\hrule\vspace{3pt}
  \newcommand{\fn}[1]{\textsc{#1}}
  \newcommand{\I}[1]{\hspace*{#1.4em}}
  \newcommand{\cm}[1]{\quad\textcolor{gray}{$\triangleright$ \textit{#1}}}
  {\small
  \begin{tabular}{@{}r@{\;\;}l@{}}
     1: & \textbf{function} \fn{3d\_segment\_mapper}$(\mathcal{P}, \mathcal{S}, k_n, T_n)$ \\
     2: & \I1$\tilde{\mathcal{S}}_n \gets \fn{segment\_keyframe}(k_n)$ \cm{2D masks (SAM)} \\
     3: & \I1$\tilde{\mathcal{P}}_n \gets \fn{project\_point\_cloud}(\mathcal{P}, T_n)$ \cm{projected map points} \\
     4: & \I1\textbf{for each} $(s, l_s) \in \tilde{\mathcal{S}}_n$ \textbf{do} \\
     5: & \I2$\mathit{mode}, v \gets \fn{get\_label\_mode\_and\_votes}(\tilde{\mathcal{P}}_n, s, \epsilon)$ \\
     6: & \I2\textbf{if} $v > \epsilon$ \textbf{then} \cm{enough evidence} \\
     7: & \I3\textbf{if} $\mathit{mode} = -1$ \textbf{then} \cm{new object} \\
     8: & \I4$S_{q+1} \gets \fn{new\_3d\_segment}(q{+}1, n, s)$;\ \ $\mathcal{S} \gets \mathcal{S} \cup \{S_{q+1}\}$;\ \ $l_s \gets q{+}1$ \\
     9: & \I3\textbf{else} \cm{already-seen instance} \\
    10: & \I4$\mathcal{S} \gets \fn{update\_3d\_segment}(S_{\mathit{mode}}, n, s)$;\ \ $l_s \gets z_l$ \\
    11: & \I1$\tilde{\mathcal{S}}_n \gets \fn{merge\_and\_prune\_2d\_segments}(\tilde{\mathcal{S}}_n)$ \\
    12: & \I1$\mathcal{P} \gets \fn{update\_pcd\_labels}(\mathcal{P}, \tilde{\mathcal{P}}_n, \tilde{\mathcal{S}}_n)$ \\
    13: & \I1\textbf{return} $\mathcal{P}, \mathcal{S}, \tilde{\mathcal{S}}_n$ \\
  \end{tabular}}
  \vspace{-2pt}\hrule
\end{algorithm}

Este bucle es la \textbf{asociación temporal} de OVO: enlaza máscaras 2D con instancias 3D
mediante proyección y voto mayoritario, sin seguimiento en 2D. Dos propiedades resultan
determinantes en el análisis posterior: la asociación se resuelve sobre la etiqueta \emph{por
punto} $l_p$, y depende por completo de que la geometría proyectada sea correcta
(Capítulos~\ref{ch:failure_modes} y siguientes).

\subsection{Instance descriptors}
\label{sec:ovo_descriptors}

El segment mapper crea y mantiene las instancias $\mathcal{S}$, pero las deja \emph{sin describir}:
el descriptor $\mathbf{d}$ de $S = (\mathbf{d}, \kappa)$ queda pendiente. Esta pieza no crea
instancias, les da semántica. Ese descriptor es un vector CLIP en el espacio lenguaje-imagen que
permite consultar la instancia con texto libre. Calcularlo es el paso más costoso de OVO, así que no
ocurre en cada keyframe: los pares $(k_n, \tilde{\mathcal{S}}_n)$ que produce la asociación esperan en la cola
$\mathcal{Q}$ (Sección~\ref{sec:ovo_segment_mapper}) y se procesan cuando hay recursos. Al sacar un
par de la cola, se calcula un descriptor por cada máscara cuyo keyframe siga vigente en el heap
$\kappa$ de su instancia, y el descriptor de la instancia se re-selecciona como el más
representativo de ese heap: la vista con menor distancia agregada al resto. La semántica de una
instancia se refina así a medida que $\kappa$ acumula mejores vistas, no en un instante fijo.

El descriptor de cada máscara 2D no es un CLIP directo, sino la fusión de tres: el del keyframe
completo, el de la máscara con el fondo eliminado y el del recorte rectangular (bounding box) que la
contiene. OVO los combina en uno mediante una media ponderada cuyos pesos \emph{por dimensión}
predice una pequeña red, en lugar de fijarlos a mano (Figura~\ref{fig:ovo_overview}): los tres pasan
por un encoder transformer y una MLP que, tras una softmax, entrega los pesos de la mezcla, una red
preentrenada con una pérdida SigLIP~\parencite{zhai2023siglip}. Las consultas se resuelven por
similitud coseno entre el texto codificado y los descriptores de las instancias
(Figura~\ref{fig:ovo_queries}).

% FIGURA: Fig. 2 del paper (pipeline del descriptor). images/ovo_overview.png
\begin{figure}[h!]
  \centering
  \includegraphics[width=\textwidth]{ovo_overview}
  \caption{The OVO pipeline for obtaining a merged CLIP descriptor per instance. Adapted
    from~\cite{berrielmartins2024ovo}.}
  \label{fig:ovo_overview}
\end{figure}

\subsection{Instance fusion at loop closure}
\label{sec:ovo_fusion_lc}

Vimos en la reconstrucción geométrica (Sección~\ref{sec:ovo_geometric}) que, tras un loop closure,
los duplicados que la deriva había separado quedan superpuestos sobre el mismo objeto. La
\emph{fusión de instancias en loop closure} es el mecanismo que OVO lanza en ese instante para
deshacerlos (Figura~\ref{fig:ovo_loop_closure}), y el punto de partida de esta tesis.

Antes de comparar nada, OVO deja el mapa a punto. Re-transforma la geometría con las poses
corregidas, como ya vimos (Sección~\ref{sec:ovo_geometric}): la nube local de cada keyframe pasa por
$\mathcal{P}_n := T'_n T_n^{-1} \mathcal{P}_n$. Los keyframes que la optimización elimina se descartan
con sus puntos, la cola $\mathcal{Q}$ de descriptores pendientes (Sección~\ref{sec:ovo_descriptors})
se vacía y se podan de $\mathcal{S}$ las instancias cuyos puntos desaparecieron por completo. Solo
entonces, sobre el mapa ya corregido, empieza la fusión.

El mecanismo de fusión existe para reunir esos duplicados superpuestos: recorre todos los pares de
instancias que quedan y decide, par a par, si son el mismo objeto registrado dos veces. Tres
condiciones han de cumplirse para fundir un par, y la implementación las comprueba en cascada, de la
más barata a la más cara, descartando el par en cuanto una falla:

\begin{itemize}
  \item \textbf{Centroide.} Los centroides de las dos nubes deben distar menos de $150\,\text{cm}$.
    Puramente geométrica y barata, elimina de un plumazo la mayoría de pares lejanos.
  \item \textbf{Descriptor.} La similitud coseno entre sus descriptores CLIP debe pasar de $0.8$,
    comprobando que comparten semántica antes de pagar el criterio más costoso.
  \item \textbf{Solapamiento.} Mide qué fracción de los puntos de la instancia más pequeña cae a
    menos de $10\,\text{cm}$ de la otra. No es un umbral único: el par pasa si esa fracción supera el
    $50\,\%$, o si supera el $20\,\%$ cuando la semántica es muy fuerte (coseno $> 0.9$), que rescata
    solapes más débiles. Es el criterio más caro, porque recorre los puntos de ambas nubes, y el
    único que comprueba que las dos ocupan de verdad el mismo volumen.
\end{itemize}

Las tres se combinan como una conjunción lógica, con umbrales fijos y comunes a toda la escena: un
par se funde solo si las cumple las tres. Fundir $S_i, S_j$ se reduce entonces a reasignar: se unen
sus listas de keyframes de referencia, $\kappa_i := \kappa_i \cup \kappa_j$, y todos los puntos
etiquetados como $j$ pasan a $i$ ($l_p := i$ para todo $P$ con $l_p = j$). La instancia $S_j$
desaparece; sus puntos y sus vistas quedan absorbidos por $S_i$, cuyo descriptor se re-selecciona
sobre el conjunto de vistas ampliado (Sección~\ref{sec:ovo_descriptors}).

Esta cascada es el único mecanismo con el que OVO une instancias, y solo contempla un escenario: el
de los duplicados que un loop closure deja superpuestos. Es, además, el punto sobre el que se
articula el resto de la tesis. Dos rasgos de su diseño lo anticipan: la
decisión es \emph{una única respuesta binaria} construida por una conjunción de compuertas, y se
toma \emph{a nivel de instancia completa}, fundiendo o no un par entero. El
Capítulo~\ref{ch:failure_modes} examina qué implica que un mecanismo con esta forma tenga que
arbitrar situaciones de naturaleza muy distinta.

% FIGURA: Fig. 5 del paper. Antes/después de loop closure. images/ovo_loop_closure.png
\begin{figure}[h!]
  \centering
  \includegraphics[width=0.72\textwidth]{ovo_loop_closure}
  \caption{Instance fusion at loop closure: drift-induced duplicates (top) merged after the loop
    closes (bottom). OVO-ORB-SLAM2, ScanNet. Adapted from~\cite{berrielmartins2024ovo}.}
  \label{fig:ovo_loop_closure}
\end{figure}

% PÁRRAFO APARCADO (2026-07-29): redundante tras el itemize geométrico/semántico.
% Aporta solo l_p (-1->relleno), el store y el diagrama. Reconciliar cuando se decida
% qué hacer con la Figura de arquitectura.
\iffalse
La Figura~\ref{fig:ovo_architecture} presenta esta arquitectura desde el ángulo que adopta esta
tesis: OVO como un \textbf{orquestador} de dos subsistemas que corren sobre el mismo flujo RGB-D
y escriben un \emph{estado compartido}. El \textbf{geometric tracking} construye y mantiene
coherente la nube de puntos densa, retroproyectando la profundidad con las poses del SLAM; el
\textbf{semantic tracking} etiqueta esos puntos en instancias y les asigna un descriptor. El punto
de contacto entre ambos es deliberadamente estrecho: cada punto de la nube lleva una etiqueta de
instancia $l_p$ que la geometría deja \emph{vacía} ($l_p = -1$) al crear el punto, y que la
semántica rellena después; el descriptor, en cambio, no vive en la nube sino en un \emph{store} de
instancias, uno por objeto. Los modelos que OVO enchufa (el backbone de SLAM, SAM y el modelo de
descriptor) quedan fuera de esa frontera, como entradas intercambiables.
\fi

\section{Analyzing the semantics for improvement}
\label{sec:ovo_factors}

Antes de señalar qué puede mejorarse en la parte semántica de OVO conviene fijar con qué se mide su
calidad, porque no todo se mide igual. Un mapa semántico puede fallar de dos formas independientes:
acertar la geometría de una instancia pero errar su etiqueta, o etiquetar bien un objeto que se
segmentó mal (partido en dos, fundido con el vecino, incompleto). Cada fallo lo capta una familia de
métricas distinta:

\begin{itemize}
  \item \textbf{Class-aware} (mIoU, mAcc). Comparan la etiqueta de cada punto con la del GT. Miden
    la calidad \emph{semántica}, que depende sobre todo del descriptor.
  \item \textbf{Class-agnostic} (AP agnóstica). Comparan la segmentación en instancias con la del GT
    \emph{sin mirar la clase}. Miden si el objeto se recuperó como una instancia correcta.
\end{itemize}

Ambas familias se definen en detalle en el Capítulo~\ref{ch:benchmark}; aquí basta la distinción,
porque es la regla con la que clasificar todo lo que sigue: cada factor que interviene golpea uno de
los dos ejes, o ambos.

Con esa regla, y el sistema ya presentado por completo, podemos dar el paso de analizarlo pieza a
pieza. El diagrama (Figura~\ref{fig:ovo_architecture}), al separar responsabilidades, aísla los
factores que condicionan la calidad semántica. Unos son \textbf{entradas externas} que OVO consume
pero no gobierna; otros, decisiones de su \textbf{propio diseño}. Los dos apartados siguientes los
recorren en ese orden, cada uno con su tabla-resumen: qué aporta cada factor, cómo falla y qué eje de
calidad compromete.

% FIGURA: arquitectura propia (dos trackings + estado compartido). images/ovo_architecture.pdf
\begin{figure}[h!]
  \centering
  \makebox[\textwidth][c]{\includegraphics[width=1.18\textwidth]{ovo_architecture}}
  \caption{OVO as an orchestrator of two trackings over a shared state.}
  \label{fig:ovo_architecture}
\end{figure}

\subsection{What OVO depends on}
\label{sec:ovo_external}

Estos factores son entradas que OVO consume pero no controla; su calidad marca el techo de lo que la
parte semántica puede lograr (Tabla~\ref{tab:factors_external}).

\begin{table}[t]
  \centering
  \small
  \setlength{\tabcolsep}{10pt}
  \renewcommand{\arraystretch}{1.15}
  \begin{tabular}{@{}l >{\raggedright\arraybackslash}p{3.7cm} >{\raggedright\arraybackslash}p{6.3cm} l@{}}
    \toprule
    \textbf{Factor} & \textbf{Rol} & \textbf{Cómo falla} & \textbf{Eje} \\
    \midrule
    Geometría & Nube por reproyección & Deriva: reproyecciones falsas, duplicados, máscaras propagadas a instancias ajenas & Agnostic \\
    \cmidrule(l{0pt}r{0pt}){1-4}
    \multirow{3}{*}{Segmentador}
      & Oportunidades de tracking & Pocas o ninguna máscara: se pierden ocasiones de crear o seguir instancias & Agnostic \\
      & Calidad de la máscara     & Una máscara errónea en mal momento siembra un error que se propaga & Agnostic \\
      & Granularidad              & Decide qué es una instancia; puede no coincidir con el GT (parte o funde objetos) & Agnostic \\
    \cmidrule(l{0pt}r{0pt}){1-4}
    \multirow{2}{*}{Descriptor}
      & Semántica de la instancia & Poco discriminativo: etiqueta pobre & Aware \\
      & Árbitro de la fusión      & Poco discriminativo: fusiones ambiguas & Both \\
    \bottomrule
  \end{tabular}
  \caption{External factors OVO consumes: their roles, how each fails, and which metric axis it hits.}
  \label{tab:factors_external}
\end{table}

La \textbf{geometría} abre la cadena. OVO no razona sobre imágenes, sino sobre la nube que
reproyecta al plano de cada keyframe, de modo que toda la asociación descansa en que esa
reproyección caiga donde debe. Con geometría ideal (poses ground-truth) el problema queda limpio y
los fallos que se observen son intrínsecos a la semántica; con un SLAM real, la deriva desplaza los
puntos, las reproyecciones mienten y aparecen duplicados y máscaras que se propagan a instancias que
no son. Por eso la fuente geométrica no es un detalle de montaje: decide si estudiamos los problemas
propios de OVO o los que le hereda el tracking.

Y hay un matiz que agrava esa dependencia. La única corrección geométrica que OVO escucha es la del
loop closure: cuando el SLAM cierra un loop, re-transforma la nube y vuelve a sincronizarla con la
trayectoria corregida. Pero el back-end refina el mapa de forma continua con bundle adjustment local,
y eso OVO no lo recoge. Entre cierres de loop, o si el loop nunca llega a producirse, la nube que
mantiene OVO se aparta poco a poco del mapa real que el SLAM sí actualiza; y como toda la asociación
se apoya en reproyecciones, esa divergencia silenciosa se traduce en errores.

El \textbf{segmentador} fija la granularidad, porque es quien decide qué es una instancia, e
interviene de tres maneras. Marca las \emph{oportunidades} de trackear: cada máscara es una ocasión
de crear o seguir una instancia, y su ausencia una ocasión perdida. Impone un techo de
\emph{calidad}: una máscara errónea en el momento equivocado siembra una instancia que ya nace mal. Y
su \emph{granularidad} puede no coincidir con la del GT, partiendo un objeto en varias instancias o
fundiendo dos en una.

El \textbf{descriptor} juega dos papeles. Fija la calidad class-aware, la etiqueta semántica de cada
instancia; y actúa como árbitro de la fusión, donde cuanto menos discriminativo sea, más difícil
resulta distinguir dos objetos parecidos de dos vistas del mismo.

\subsection{What OVO decides}
\label{sec:ovo_internal}

Frente a esas entradas, estos dos factores son decisiones internas de OVO, y por tanto los que la
tesis puede intervenir directamente (Tabla~\ref{tab:factors_internal}).

\begin{table}[h!]
  \centering
  \small
  \setlength{\tabcolsep}{10pt}
  \renewcommand{\arraystretch}{1.15}
  \begin{tabular}{@{}l >{\raggedright\arraybackslash}p{3.7cm} >{\raggedright\arraybackslash}p{6.3cm} l@{}}
    \toprule
    \textbf{Factor} & \textbf{Rol} & \textbf{Cómo falla} & \textbf{Eje} \\
    \midrule
    Creación inicial & Alta de instancias & Decisión inmutable: el error inicial no se deshace, se propaga & Agnostic \\
    \cmidrule(l{0pt}r{0pt}){1-4}
    Fusión & Corrección de duplicados & Solo une solapados; ciega a splits y a sobre-segmentación & Both \\
    \bottomrule
  \end{tabular}
  \caption{OVO's own design decisions: their roles, how each fails, and which metric axis it hits.}
  \label{tab:factors_internal}
\end{table}

La \textbf{creación inicial} de instancias es imperativa e \textbf{inmutable}: la asociación
(Sección~\ref{sec:ovo_segment_mapper}) da de alta la instancia en su primer avistamiento y no hay
vuelta atrás. Un error de segmentación en ese instante no se corrige; se propaga a cada keyframe
siguiente por el propio mecanismo de propagación.

La \textbf{fusión} es el único mecanismo correctivo de OVO, y su alcance es estrecho: solo actúa
cuando dos instancias se \emph{solapan} tras un loop closure. Esa condición, el solapamiento, deja
fuera clases enteras de problemas. Una instancia que nació mal segmentada no se solapa con nada que
la arregle; una máscara que a lo largo de la secuencia se disputan varias instancias no produce un
solapamiento que la fusión pueda arbitrar; y una granularidad que no casa con la del GT, un objeto
partido en trozos o dos fundidos en uno, tampoco se corrige, porque no hay dos nubes superpuestas que
unir. En ninguno de esos casos existe el overlap del que la fusión depende, de modo que OVO,
sencillamente, carece de un mecanismo que los aborde: los fallos, una vez sembrados, no tienen quién
los repare (Sección~\ref{sec:ovo_limitations}).

\section{Limitations}
\label{sec:ovo_limitations}

Los propios autores de OVO señalan tres limitaciones del sistema~\parencite{berrielmartins2024ovo}:
la generalización a escenas outdoor a gran escala (con distribuciones de clase, iluminación y errores
de tracking distintos), un ligero sesgo del CLIP merging hacia las clases vistas en entrenamiento, y
que la fusión de instancias en loop closure, aunque corrige eficazmente la deriva odométrica,
\enquote{a veces no llega a fundir instancias que sí deberían fundirse}. Las dos primeras caen fuera
del alcance de este trabajo; es la tercera la que tomamos como punto de partida y situamos en un
marco más amplio: el mecanismo de fusión arrastra tres limitaciones estructurales.

\begin{itemize}
  \item \textbf{Una sola decisión binaria para varias preguntas distintas.} Aceptar o rechazar un
    par depende, a la vez, de si la geometría está desactualizada, de si la segmentación 2D venía
    mal, de si el descriptor es ambiguo y de si el objeto es uno o son dos. Un único umbral no
    puede arbitrar cuatro preguntas independientes.
  \item \textbf{Granularidad todo-o-nada.} La decisión se toma sobre la instancia completa, pero
    muchos fallos son \emph{parciales} (parte de un objeto bien, parte mal). Un fallo parcial solo
    se puede corregir con una decisión parcial, a nivel de punto.
  \item \textbf{Se presupone que la geometría sigue siendo válida.} La fusión decide sobre las
    posiciones actuales de los puntos, pero el Bundle Adjustment local
    sigue deformando el mapa por debajo, sin que la capa semántica se entere.
\end{itemize}

Estas tres limitaciones organizan el resto del trabajo: primero hace falta poder \emph{medir} una
decisión de fusión (Capítulo~\ref{ch:benchmark}); después, diagnosticar por qué falla
(Capítulo~\ref{ch:failure_modes}); y a partir de ahí atacar cada limitación por separado.
